\section{Background and Related Work}

Structured decision models replace free-form generation with a typed question
and a typed answer. Jev, the commercial System One model from TypeSafe, accepts
a state and a set of questions and returns \noul{} (a binary judgment),
\choice{} (a selection among named candidates) or \score{} (an ordered rating)
with associated probabilities \citep{typesafe2026}. \citet{deusser2026jev}
benchmark Jev on 37 datasets, and \citet{guo2026justask} apply it to alignment
failure detection; both treat the model as a general decision interface rather
than a text generator. TypeSafe describes the training objective as
Reinforcement Learning for Calibrated Decisions (RLCD), which optimises the
reported distribution with strictly proper scoring rules so that probabilities
carry information beyond the argmax; the full recipe is not public.
\citet{gao2026openjev} provide an independent open implementation. The name
RLCD also appears in \citet{yang2023rlcd} for Reinforcement Learning from
Contrast Distillation; the two are unrelated, and we use RLCD only in the
calibrated-decision sense.

The two models we run differ in backbone and supervision while sharing the
interface. Laya 0.4.16 is a ModernBERT-large encoder with decision-specific
heads, about 421M parameters, and supports \noul{}, \choice{} and \score{} in a
non-autoregressive pass; it is the open model closest to the RLCD description
\citep{laya2026,layastudio2026}. Vela 2.0 0.3B is a 22-layer ModernBERT encoder
with an 8{,}192-token window; it exposes the same three actions plus \spanact{}
and \setact{}, and was initialised from the Decision-1.0-Kai \choice{} path with
task-specific losses and routing data \citep{vllm2026vela,warner2024modernbert}.
Vela's card does not call the recipe RLCD, so we treat Laya as an open RLCD
model and Vela as an interface-matched comparison with a different training
route (Appendix~\ref{app:models}).

Calibration is the natural first lens on these probabilities. Neural networks
are often overconfident, and temperature scaling or other post-hoc maps can
reduce expected calibration error without changing the argmax
\citep{guo2017calibration}. Language models can also partially recognise their
own errors, which makes confidence a plausible routing feature
\citep{kadavath2022know}, and proper scoring rules formalise why a reported
probability should be rewarded for honesty rather than sharpness alone
\citep{gneiting2007scoring}. All of this work concerns the probability of the
answer the model has already produced, not whether that probability orders the
counterfactual value of a different question type. We measure ECE and Brier for
current-answer correctness, then measure the same probability's ability to rank
rows that the alternative action would repair.

Cascades and routers allocate computation across models or actions. The
classical cascade rejects easy examples early and sends hard ones to a stronger
classifier \citep{viola2001rapid}; modern systems learn when to call a larger
language model \citep{chen2024frugalgpt,aggarwal2024automix,ong2024routellm},
how many candidates to sample \citep{ding2025bestroute}, or how to combine
routing and cascading in one objective \citep{dekoninck2025unified}. The
theoretical caution comes from \citet{jitkrittum2023cascade}: confidence-based
deferral is optimal only when downstream errors are sufficiently tied to
upstream confidence, and it can fail under specialised downstream errors, label
noise or distribution shift. \citet{damani2025hard} argue that computation should
follow expected benefit rather than surface difficulty. Our experiments
instantiate this warning in a structured decision model: the second action is a
different question to the same frozen network, not a larger model, and its value
depends on an error structure that the first probability does not encode. For
ProcessBench, selective prediction supplies the abstention framing:
\citet{chow1970optimum} derives the optimal reject rule under known
distributions, and \citet{geifman2017selective}, \citet{kamath2020selective} and
\citet{xin2021art} study selective classification, question answering and the
separation of coverage from conditional accuracy.

Evolutionary methods are a reasonable choice when the policy is a small
non-differentiable object and the objective is a cached replay. Genetic
algorithms, evolution strategies and NSGA-II cover single-objective and Pareto
formulations \citep{deb2002nsga2,hansen2001cmaes,stanley2002neat}, and
regularised evolution has been used for architecture search
\citep{real2019regularized,elsken2019nas}. In NLP, EvoPrompt and Promptbreeder
evolve prompts or mutation instructions with language models in the loop
\citep{guo2024evoprompt,fernando2024promptbreeder}. Those systems search over
text that changes what the model sees; we hold the model and the question
template fixed and search over a two-parameter gate that decides whether a
second typed question is asked. Prompt search can improve a single call, whereas
our search can only exploit complementarity that already exists between two
fixed actions. A flat landscape is therefore a finding about the action set, not
merely about optimiser strength.

Our protocol follows the variance-aware reporting practices of
\citet{henderson2018matters}, \citet{bouthillier2021variance} and
\citet{dodge2019show}: we fix the action set, templates and budget before
looking at confirm, report every seed-level interval rather than only the best
cell, and label the locked test as not run rather than substituting a confirm
number for it.
